% !TeX program = pdflatex
\documentclass[11pt,a4paper]{article}

\usepackage[T1]{fontenc}
\usepackage[utf8]{inputenc}
\IfFileExists{libertinus.sty}{\usepackage{libertinus}}{\usepackage{lmodern}}
\usepackage{microtype}
\usepackage[a4paper,margin=27mm,headheight=15pt]{geometry}
\usepackage{amsmath,amssymb,amsthm,mathtools}
\usepackage{bm}
\usepackage{mathrsfs}
\usepackage{booktabs,longtable,array,tabularx}
\usepackage{enumitem}
\usepackage{xcolor}
\usepackage{tikz}
\usetikzlibrary{arrows.meta,positioning,calc,decorations.pathreplacing}
\usepackage[most]{tcolorbox}
\usepackage{fancyhdr}
\usepackage{xurl}
\usepackage{listings}
\usepackage{hyperref}
\usepackage[nameinlink,noabbrev,capitalise]{cleveref}

\definecolor{deepolive}{RGB}{67,79,45}
\definecolor{softolive}{RGB}{236,239,226}
\definecolor{darkteal}{RGB}{31,83,84}
\definecolor{warmgray}{RGB}{92,87,78}
\definecolor{softcoral}{RGB}{246,225,219}
\definecolor{linkblue}{RGB}{31,76,122}
\definecolor{codegray}{RGB}{246,246,244}

\hypersetup{
  colorlinks=true,
  linkcolor=deepolive,
  citecolor=darkteal,
  urlcolor=linkblue,
  pdftitle={The Complete Asymptotic Transseries of the Fubini Numbers},
  pdfauthor={A self-contained analytic-combinatorial synthesis},
  pdfsubject={Ordered Bell numbers, exact pole expansions, transseries, Bell-polynomial coefficient calculus, and resurgence}
}

\setlist{itemsep=2pt,topsep=4pt,parsep=1pt}
\allowdisplaybreaks[3]
\raggedbottom
\emergencystretch=2em
\setcounter{tocdepth}{2}
\setcounter{secnumdepth}{3}

\pagestyle{fancy}
\fancyhf{}
\fancyhead[L]{\small\nouppercase{\leftmark}}
\fancyhead[R]{\small\thepage}
\renewcommand{\headrulewidth}{0.4pt}

\newtheorem{theorem}{Theorem}[section]
\newtheorem{lemma}[theorem]{Lemma}
\newtheorem{proposition}[theorem]{Proposition}
\newtheorem{corollary}[theorem]{Corollary}
\newtheorem{identity}[theorem]{Identity}
\theoremstyle{definition}
\newtheorem{definition}[theorem]{Definition}
\newtheorem{example}[theorem]{Example}
\newtheorem{algorithm}[theorem]{Algorithm}
\theoremstyle{remark}
\newtheorem{remark}[theorem]{Remark}

\newtcolorbox{masterbox}[1][]{
  enhanced,breakable,colback=softolive,colframe=deepolive,
  boxrule=0.7pt,arc=1mm,left=2mm,right=2mm,top=1.5mm,bottom=1.5mm,
  title={Master result},fonttitle=\bfseries,#1
}
\newtcolorbox{insightbox}[1][]{
  enhanced,breakable,colback=softcoral,colframe=warmgray,
  boxrule=0.6pt,arc=1mm,left=2mm,right=2mm,top=1.5mm,bottom=1.5mm,
  title={Structural point},fonttitle=\bfseries,#1
}
\newtcolorbox{notationbox}[1][]{
  enhanced,breakable,colback=white,colframe=darkteal,
  boxrule=0.6pt,arc=1mm,left=2mm,right=2mm,top=1.5mm,bottom=1.5mm,
  title={Notation},fonttitle=\bfseries,#1
}

\lstset{
  basicstyle=\ttfamily\small,
  backgroundcolor=\color{codegray},
  frame=single,
  rulecolor=\color{warmgray},
  breaklines=true,
  columns=fullflexible,
  keepspaces=true,
  showstringspaces=false,
  aboveskip=8pt,
  belowskip=8pt
}

\newcommand{\Fub}{\mathsf F}
\newcommand{\Stirling}[2]{\left\{\!\!\begin{matrix}#1\\#2\end{matrix}\!\!\right\}}
\newcommand{\e}{\mathrm e}
\newcommand{\ii}{\mathrm i}
\newcommand{\dd}{\,\mathrm d}
\newcommand{\Res}{\operatorname*{Res}}
\newcommand{\Log}{\operatorname{Log}}
\newcommand{\Li}{\operatorname{Li}}
\newcommand{\Borel}{\mathcal B}
\newcommand{\Laplace}{\mathcal L}
\newcommand{\cO}{\mathcal O}
\newcommand{\abs}[1]{\left|#1\right|}
\newcommand{\norm}[1]{\left\lVert#1\right\rVert}
\newcommand{\falling}[2]{#1^{\underline{#2}}}
\newcommand{\rising}[2]{#1^{\overline{#2}}}
\newcommand{\coeff}[2]{\left[#1^{#2}\right]}
\newcommand{\N}{\mathbb N}
\newcommand{\Z}{\mathbb Z}
\newcommand{\C}{\mathbb C}
\newcommand{\R}{\mathbb R}

\title{\Huge\bfseries The Complete Asymptotic Transseries\\[2mm]
       of the Fubini Numbers\\[4mm]
       \Large Exact Pole Sectors, All-Orders Coefficients,\\
       Certified Remainders, and a Resurgent Interpretation}
\author{A self-contained analytic-combinatorial synthesis}
\date{September 2026}

\begin{document}
\maketitle

\begin{abstract}
The Fubini numbers, also called ordered Bell numbers or preferential-arrangement
numbers, have exponential generating function
\(
 (2-\e^z)^{-1}
\).
Its singularities form the vertical lattice
\(
 \rho_k=\log 2+2\pi\ii k
\), \(k\in\Z\), and every singularity is a simple pole of residue
\(-1/2\).  This elementary meromorphic geometry makes it possible to go much
further than the familiar leading estimate
\(
 \Fub_n\sim n!/[2(\log 2)^{n+1}]
\).
For every \(n\ge1\) we prove the exact, absolutely convergent identity
\[
 \Fub_n=\frac{n!}{2}\sum_{k\in\Z}
          (\log 2+2\pi\ii k)^{-n-1}.
\]
After normalization by the dominant pole, this is already the complete
exponential--oscillatory transseries: every nonzero pole gives one sector, no
sector carries an inverse-power descendant tower, and there are no hidden
Fubini-specific sectors.  Pairing conjugate poles gives a real expansion whose
amplitudes, actions, and phases are explicit.

We derive explicit and elementary truncation bounds in terms of the Hurwitz zeta
function, obtain certified nearest-integer reconstruction criteria, and explain
why exponentially small relative errors may nevertheless be enormous in
absolute value.  Expanding the exact factorial prefactor by Stirling's series
produces the conventional two-level transseries.  Its coefficient in pole
sector \(k\) and algebraic order \(j\) is
\[
 C_{k,j}=\frac{\log 2}{\log 2+2\pi\ii k}\,\gamma_j,
\]
where \(\gamma_j\) is given explicitly by complete exponential Bell
polynomials, by a partition sum, and by a Bernoulli-number recurrence.  A
multi-index coefficient calculus is then developed for arbitrary powers,
logarithms, reciprocals, and other analytic functions of the normalized
sequence.  Finally, the formal ordinary generating series is shown to have
Borel transform exactly \((2-\e^\xi)^{-1}\), so the pole lattice is at once the
Borel singularity set and the complete large-order data.  Extensions to Fubini
polynomials and to higher powers \((2-\e^z)^{-r}\) are included, the latter
having general coefficients expressed by N\"orlund (generalized Bernoulli)
numbers.
\end{abstract}

\clearpage
\begin{masterbox}
Put
\[
 \alpha=\log 2,\qquad \rho_k=\alpha+2\pi\ii k,\qquad
 \omega_k=\frac{\alpha}{\rho_k}.
\]
Then, for every integer \(n\ge1\),
\[
 \boxed{\displaystyle
 \frac{2\alpha^{n+1}}{n!}\Fub_n
 =\sum_{k\in\Z}\omega_k^{n+1}}
\]
is an \emph{exact} identity, not merely an asymptotic expansion.  Since
\(\omega_0=1\) and \(\abs{\omega_k}<1\) for \(k\ne0\), it is simultaneously
the complete factorial-normalized transseries.
\end{masterbox}

\tableofcontents
\newpage

\section{Introduction and scope}

\subsection{From ordered partitions to a meromorphic generating function}

A partition of a labelled \(n\)-element set becomes an \emph{ordered set
partition} when its blocks are arranged in a linear order.  If there are \(j\)
blocks, then the underlying partition can be chosen in
\(\Stirling{n}{j}\) ways and the blocks can be ordered in \(j!\) ways.
Consequently the \(n\)th Fubini number is
\begin{equation}\label{eq:def-fubini}
 \Fub_n=\sum_{j=0}^n j!\Stirling{n}{j},
 \qquad n\ge0.
\end{equation}
The first values are
\[
 1,1,3,13,75,541,4683,47293,545835,7087261,\ldots.
\]
The same sequence counts weak orderings, preferential arrangements, and a
variety of equivalent structures; see Comtet~\cite{Comtet},
Bailey~\cite{Bailey}, Pippenger~\cite{Pippenger}, and the references in
OEIS A000670~\cite{OEIS}.

The labelled sequence construction immediately gives
\begin{equation}\label{eq:egf}
 G(z):=\sum_{n\ge0}\Fub_n\frac{z^n}{n!}
 =\sum_{j\ge0}(\e^z-1)^j
 =\frac{1}{2-\e^z}.
\end{equation}
The linked \emph{Combinatorial Coefficient Calculus} develops precisely this
Stirling-number and exponential-generating-function foundation
\cite{ReshetnikovCCC}.  The linked series-and-transseries collection develops
the formal operations used after an asymptotic scale has been identified
\cite{ReshetnikovTransseries}.  The present article joins those two viewpoints:
it extracts the complete transseries from the meromorphic geometry of
\eqref{eq:egf}, and then expresses every coefficient in the Bell-polynomial
language of coefficient calculus.

\subsection{Why an ordinary Poincar\'e series is not enough}

The nearest singularity of \(G\) to the origin is the positive real pole
\(\alpha=\log 2\), so elementary singularity analysis
\cite[Chap.~IV]{FlajoletSedgewick} yields
\begin{equation}\label{eq:leading}
 \Fub_n\sim \frac{n!}{2\alpha^{n+1}}.
\end{equation}
This classical equivalent goes back at least to Barth\'elemy's analysis of
total preorders~\cite{Barthelemy}.  It suppresses, however, an infinite lattice
of further poles at \(\alpha+2\pi\ii k\).  Their contributions are smaller than
every power of \(1/n\) relative to the dominant term, and therefore an ordinary
Poincar\'e series sees only
\[
 \frac{2\alpha^{n+1}\Fub_n}{n!}
 \sim 1+\frac0n+\frac0{n^2}+\cdots.
\]
The zero coefficients are correct but incomplete: the actual error is a
structured sum of exponentially small, oscillatory terms.  A transseries keeps
those terms rather than discarding them.

\subsection{What ``full'' means in this article}

There are two natural normalizations, and it is useful to distinguish them.

\begin{enumerate}
\item \textbf{Factorial-normalized expansion.}
      Keep \(n!\) exact.  The pole sum in the abstract is then exact and contains
      every Fubini-specific exponential sector.  Each simple pole contributes
      one transmonomial and no inverse-power descendants.
\item \textbf{Elementary Stirling normalization.}
      Expand \(n!\) in powers of \(1/n\).  Every pole sector then inherits the
      same universal Stirling tower.  We derive its coefficients to arbitrary
      order and give a closed formula for the coefficient attached to every
      pair \((k,j)\) of pole index and algebraic order.
\end{enumerate}

The first form is analytically stronger because it is exact.  The second is the
form usually expected when a transseries is written over the elementary scale
\(
 \sqrt{n}(n/(\e\alpha))^n
\).
Together they leave no unspecified Fubini-specific remainder.

\section{Combinatorial and exact analytic foundations}

\subsection{The exponential generating function and recurrence}

\begin{theorem}[Basic Fubini identities]\label{thm:basic}
For \(n\ge0\), the numbers in \eqref{eq:def-fubini} satisfy
\begin{align}
 \sum_{n\ge0}\Fub_n\frac{z^n}{n!}&=\frac1{2-\e^z},
 \label{eq:basic-egf}\\
 \Fub_0&=1,\qquad
 \Fub_n=\sum_{j=1}^n\binom nj\Fub_{n-j}\quad(n\ge1),
 \label{eq:basic-rec}\\
 \Fub_n&=\frac12\sum_{m=0}^{\infty}\frac{m^n}{2^m}.
 \label{eq:dobinski-fubini}
\end{align}
In \eqref{eq:dobinski-fubini}, the convention \(0^0=1\) is used when \(n=0\).
For \(n\ge1\), the term \(m=0\) vanishes.
\end{theorem}

\begin{proof}
The combinatorial derivation of \eqref{eq:basic-egf} was given in
\eqref{eq:egf}.  Multiplying that identity by \(2-\e^z\) and comparing the
coefficient of \(z^n/n!\) gives
\[
 2\Fub_n-\sum_{j=0}^n\binom nj\Fub_{n-j}=0.
\]
The \(j=0\) summand is \(\Fub_n\), and moving it to the other side gives
\eqref{eq:basic-rec}.  Finally,
\[
 \frac1{2-\e^z}
 =\frac12\frac1{1-\e^z/2}
 =\frac12\sum_{m\ge0}\frac{\e^{mz}}{2^m}.
\]
Extracting the coefficient of \(z^n/n!\) proves
\eqref{eq:dobinski-fubini}.  The interchange is valid analytically near
\(z=0\), and also as an identity of formal exponential series because every
coefficient is an absolutely convergent geometric moment.
\end{proof}

\begin{remark}[Probabilistic interpretation]
Let \(X\) have the geometric distribution
\(
 \Pr(X=m)=2^{-m-1}
\) for \(m=0,1,2,\ldots\).  Then
\(
 \Fub_n=\mathbb E[X^n]
\).
Thus the same moments that count ordered partitions also measure powers of the
number of failures before the first success in fair Bernoulli trials.
\end{remark}

\begin{remark}[Polylogarithm form]
For \(n\ge1\),
\begin{equation}\label{eq:polylog}
 \Fub_n=\frac12\Li_{-n}\!\left(\frac12\right).
\end{equation}
This is exact, but the pole representation developed below is better adapted to
large \(n\).
\end{remark}

\subsection{The continuum term already gives the dominant asymptotic}

Set
\begin{equation}\label{eq:alpha}
 \alpha:=\log 2.
\end{equation}
The moment formula may be written
\begin{equation}\label{eq:moment-alpha}
 \Fub_n=\frac12\sum_{m\ge1}m^n\e^{-\alpha m},
 \qquad n\ge1.
\end{equation}
Replacing the lattice sum by its continuum integral gives
\begin{equation}\label{eq:continuum}
 \frac12\int_0^\infty x^n\e^{-\alpha x}\dd x
 =\frac{\Gamma(n+1)}{2\alpha^{n+1}}
 =\frac{n!}{2\alpha^{n+1}}.
\end{equation}
The remarkable fact is that the difference between \eqref{eq:moment-alpha}
and \eqref{eq:continuum} is not an inverse-power Euler--Maclaurin series.  For
fixed positive \(\alpha\), all endpoint derivatives that could contribute at
algebraic orders are encoded by an exact periodic pole sum, and the relative
difference is exponentially small in \(n\).  Poisson summation will make this
statement exact.

\section{The singularity lattice and the exact pole formula}

\subsection{Location and residues of all singularities}

The poles of \(G(z)=(2-\e^z)^{-1}\) solve \(\e^z=2\).  Hence
\begin{equation}\label{eq:rho}
 \rho_k:=\alpha+2\pi\ii k,
 \qquad k\in\Z.
\end{equation}
They are all simple, and
\begin{equation}\label{eq:residue}
 \Res_{z=\rho_k}G(z)
 =\frac1{-\e^{\rho_k}}=-\frac12.
\end{equation}
Thus every pole has the same residue; only its distance and phase from the
origin vary.

\begin{figure}[ht]
\centering
\begin{tikzpicture}[x=2.4cm,y=0.42cm,>=Latex]
  \draw[->] (-0.35,0)--(1.55,0) node[right] {$\Re z$};
  \draw[->] (0,-14.5)--(0,14.5) node[above] {$\Im z$};
  \draw[dashed,darkteal] (0.693147,-14)--(0.693147,14);
  \foreach \k in {-2,-1,0,1,2}{
    \fill[deepolive] (0.693147,{6.283185*\k}) circle (2.1pt);
  }
  \node[below] at (0.693147,0) {$\rho_0=\log 2$};
  \node[right] at (0.693147,6.283185) {$\rho_1$};
  \node[right] at (0.693147,-6.283185) {$\rho_{-1}$};
  \node[right] at (0.693147,12.56637) {$\rho_2$};
  \node[right] at (0.693147,-12.56637) {$\rho_{-2}$};
  \draw[decorate,decoration={brace,amplitude=4pt},warmgray]
       (0.79,0)--(0.79,6.283185)
       node[midway,right=6pt] {$2\pi$};
\end{tikzpicture}
\caption{The complete singularity lattice of \(G(z)=(2-\e^z)^{-1}\).
Every point on the line \(\Re z=\log2\) is a simple pole of residue \(-1/2\).}
\label{fig:poles}
\end{figure}

\subsection{A Mittag--Leffler identity}

The following classical partial fraction identity is the analytic engine behind
the exact expansion; see also \cite{WhittakerWatson}.

\begin{lemma}[Partial fractions for the Bose--Einstein kernel]\label{lem:coth}
For \(w\notin2\pi\ii\Z\),
\begin{equation}\label{eq:coth-pf}
 \frac1{\e^w-1}
 =\frac1w-\frac12
  +2w\sum_{k=1}^{\infty}\frac1{w^2+4\pi^2k^2}.
\end{equation}
The series converges locally uniformly away from the poles.  Equivalently,
\begin{equation}\label{eq:coth-sym}
 \frac12\coth\frac w2
 =\frac1w+2w\sum_{k=1}^{\infty}\frac1{w^2+4\pi^2k^2}.
\end{equation}
\end{lemma}

\begin{proof}
The canonical product
\[
 \frac{\sinh(w/2)}{w/2}
 =\prod_{k=1}^{\infty}
   \left(1+\frac{w^2}{4\pi^2k^2}\right)
\]
converges locally uniformly.  Taking its logarithmic derivative gives
\eqref{eq:coth-sym}.  Since
\(
 (\e^w-1)^{-1}=\tfrac12\coth(w/2)-\tfrac12
\), equation \eqref{eq:coth-pf} follows.
\end{proof}

\subsection{Exact pole expansion}

\begin{theorem}[Exact lattice formula]\label{thm:exact-poles}
For every integer \(n\ge1\),
\begin{equation}\label{eq:exact-poles}
 \boxed{\displaystyle
 \Fub_n=\frac{n!}{2}\sum_{k\in\Z}\rho_k^{-n-1}}
 \qquad
 \left(\rho_k=\log2+2\pi\ii k\right).
\end{equation}
The series is absolutely convergent.
\end{theorem}

\begin{proof}
From \eqref{eq:moment-alpha},
\begin{equation}\label{eq:derivative-moment}
 2\Fub_n
 =\sum_{m\ge1}m^n\e^{-\alpha m}
 =(-1)^n\frac{\dd^n}{\dd\alpha^n}
   \frac1{\e^\alpha-1}.
\end{equation}
Differentiate \eqref{eq:coth-pf} \(n\) times.  For \(n\ge1\), the constant
\(-1/2\) disappears, and the differentiated partial fractions can be paired as
\[
 \frac1{\alpha+2\pi\ii k}
 \quad\text{and}\quad
 \frac1{\alpha-2\pi\ii k}.
\]
Local uniform convergence justifies termwise differentiation.  Since
\[
 (-1)^n\frac{\dd^n}{\dd\alpha^n}
 \frac1{\alpha+2\pi\ii k}
 =\frac{n!}{(\alpha+2\pi\ii k)^{n+1}},
\]
substitution into \eqref{eq:derivative-moment} yields
\eqref{eq:exact-poles}.  Absolute convergence follows from
\(
 \abs{\rho_k}^{-n-1}=\cO(\abs{k}^{-n-1})
\), whose sum converges for \(n\ge1\).
\end{proof}

\begin{remark}[The exceptional index \(n=0\)]
Formula \eqref{eq:exact-poles} is stated for \(n\ge1\) because the corresponding
two-sided pole series is not absolutely convergent and requires symmetric
summation; moreover, the constant \(-1/2\) in
\eqref{eq:coth-pf} no longer disappears.  In fact, symmetric summation gives
\[
 \sum_{k\in\Z}^{\mathrm{sym}}\frac1{\alpha+2\pi\ii k}
 =\frac12\coth\!\left(\frac\alpha2\right)=\frac32,
 \qquad \alpha=\log 2,
\]
so the naively continued right-hand side of \eqref{eq:exact-poles} is
\(3/4\), not \(\Fub_0=1\).  The bookkeeping is transparent in
\eqref{eq:coth-pf}: multiplying the regularization constant \(-1/2\) by the
outer factor \(1/2\) changes \(3/4\) to the positive-\(m\) moment \(1/2\),
and the separate \(m=0\) term in \eqref{eq:dobinski-fubini} contributes the
remaining \(1/2\).  Thus the exact pole identity is naturally an
\(n\ge1\) statement.
\end{remark}

\subsection{A second proof by one-sided Poisson summation}

The pole formula also has a transparent harmonic-analysis interpretation.

\begin{proposition}[Poisson form of the lattice correction]\label{prop:poisson}
For \(n\ge1\), define
\[
 f_n(x)=\begin{cases}
 x^n\e^{-\alpha x},&x\ge0,\\
 0,&x<0.
 \end{cases}
\]
With Fourier transform
\(
 \widehat f_n(k)=\int_{\R}f_n(x)\e^{-2\pi\ii kx}\dd x
\), one has
\begin{equation}\label{eq:fourier-f}
 \widehat f_n(k)=\frac{n!}{(\alpha+2\pi\ii k)^{n+1}}.
\end{equation}
Poisson summation gives
\begin{equation}\label{eq:poisson}
 \sum_{m\ge1}m^n\e^{-\alpha m}
 =\sum_{k\in\Z}\frac{n!}{(\alpha+2\pi\ii k)^{n+1}},
\end{equation}
and hence \eqref{eq:exact-poles}.
\end{proposition}

\begin{proof}
The Laplace integral gives \eqref{eq:fourier-f}.  The function \(f_n\) is
continuous at the origin for \(n\ge1\), piecewise smooth, exponentially
decaying, and its transform is absolutely summable because it is
\(\cO(\abs{k}^{-n-1})\).  The standard one-sided form of Poisson summation is
therefore applicable.  Since \(f_n(m)=0\) for negative integers and
\(f_n(0)=0\), its left side is the sum in \eqref{eq:poisson}.  Combining this
with \eqref{eq:moment-alpha} proves the result.
\end{proof}

\begin{insightbox}
The \(k=0\) Fourier mode in \eqref{eq:poisson} is exactly the continuum integral
\eqref{eq:continuum}.  Every \(k\ne0\) mode is an exponentially small lattice
correction.  Thus the full transseries can be read either from the poles of the
exponential generating function or from the nonzero Fourier harmonics of the
geometric moment sum.
\end{insightbox}

\section{The complete factorial-normalized transseries}

\subsection{Actions, amplitudes, and phases}

Define, for \(k\in\Z\),
\begin{equation}\label{eq:omega-lambda}
 \omega_k:=\frac{\alpha}{\rho_k},
 \qquad
 \Lambda_k:=\Log\frac{\rho_k}{\alpha}=-\Log\omega_k,
\end{equation}
where \(\Log\) is the principal logarithm.  For \(k\ge1\), also put
\begin{align}
 r_k&:=\abs{\rho_k}=\sqrt{\alpha^2+4\pi^2k^2},
 \label{eq:rk}\\
 q_k&:=\abs{\omega_k}=\frac{\alpha}{r_k},
 \label{eq:qk}\\
 \lambda_k&:=\Re\Lambda_k
 =\log\frac{r_k}{\alpha}
 =\frac12\log\left(1+\frac{4\pi^2k^2}{\alpha^2}\right),
 \label{eq:lambdak}\\
 \theta_k&:=\arg\rho_k
 =\arctan\frac{2\pi k}{\alpha}.
 \label{eq:thetak}
\end{align}
Then \(q_k=\e^{-\lambda_k}\),
\(
 \omega_k=q_k\e^{-\ii\theta_k}
\), and
\(
 \omega_{-k}=\overline{\omega_k}
\).

\begin{theorem}[Complete exact transseries]\label{thm:complete-transseries}
For every \(n\ge1\),
\begin{align}
 \frac{2\alpha^{n+1}}{n!}\Fub_n
 &=\sum_{k\in\Z}\omega_k^{n+1}
 \label{eq:complex-transseries}\\
 &=1+2\sum_{k=1}^{\infty}
 q_k^{n+1}\cos\bigl((n+1)\theta_k\bigr)
 \label{eq:real-transseries}\\
 &=1+2\sum_{k=1}^{\infty}
 \e^{-(n+1)\lambda_k}
 \cos\bigl((n+1)\theta_k\bigr).
 \label{eq:action-transseries}
\end{align}
All three expressions are exact and absolutely convergent.
For every fixed \(K\ge0\),
\begin{equation}\label{eq:finite-trans-asymp}
 \frac{2\alpha^{n+1}}{n!}\Fub_n
 =1+2\sum_{k=1}^{K}q_k^{n+1}
       \cos\bigl((n+1)\theta_k\bigr)
  +\cO_K\!\left(q_{K+1}^{n+1}\right)
\end{equation}
as \(n\to\infty\).
\end{theorem}

\begin{proof}
Equation \eqref{eq:complex-transseries} is just
\eqref{eq:exact-poles} multiplied by \(2\alpha^{n+1}/n!\).
Pairing the terms indexed by \(k\) and \(-k\) gives
\eqref{eq:real-transseries}, and \(q_k=\e^{-\lambda_k}\) gives
\eqref{eq:action-transseries}.  Absolute convergence was already established.
For fixed \(K\), the omitted magnitudes form a convergent sequence dominated by
its first scale times a constant independent of sufficiently large \(n\):
\[
 \sum_{k>K}q_k^{n+1}
 =q_{K+1}^{n+1}
   \left(1+\sum_{k>K+1}
   \left(\frac{q_k}{q_{K+1}}\right)^{n+1}\right)
 =\cO_K(q_{K+1}^{n+1}).
\]
This proves \eqref{eq:finite-trans-asymp}.
\end{proof}

The first few sector parameters are displayed in \cref{tab:sectors}.

\begin{table}[ht]
\centering
\caption{The first exponential--oscillatory sectors.  Here
\(q_k=\e^{-\lambda_k}\) and \(\theta_k/\pi\) is the normalized phase.}
\label{tab:sectors}
\begin{tabular}{@{}rllll@{}}
\toprule
\(k\)&\(r_k\)&\(q_k\)&\(\lambda_k\)&\(\theta_k/\pi\)\\
\midrule
1&6.3213029209&0.1096525810&2.2104382659&0.4650261734\\
2&12.5854727139&0.0550751804&2.8990561102&0.4824601508\\
3&18.8622960281&0.0367477628&3.3036779310&0.4883001896\\
4&25.1422977208&0.0275689672&3.5910645166&0.4912234132\\
5&31.4235722527&0.0220581917&3.8140712405&0.4929780900\\
6&37.7054835106&0.0183831930&3.9963184557&0.4941481183\\
\bottomrule
\end{tabular}
\end{table}

In particular, the first exponentially improved approximation is
\begin{equation}\label{eq:first-pair}
 \Fub_n=\frac{n!}{2\alpha^{n+1}}
 \left\{1+2q_1^{n+1}
 \cos\bigl((n+1)\theta_1\bigr)
 +\cO(q_2^{n+1})\right\}.
\end{equation}
The correction oscillates because the nearest nonreal singularities occur as a
conjugate pair.

\subsection{Why every algebraic correction vanishes after factorial normalization}

\begin{corollary}[Flatness beyond all powers]\label{cor:flat}
For every fixed \(M\ge1\),
\begin{equation}\label{eq:flat}
 \frac{2\alpha^{n+1}}{n!}\Fub_n
 =1+o(n^{-M}).
\end{equation}
Equivalently, the complete Poincar\'e expansion of the normalized ratio is
\begin{equation}\label{eq:poincare-zero}
 \frac{2\alpha^{n+1}}{n!}\Fub_n
 \sim 1+\frac0n+\frac0{n^2}+\frac0{n^3}+\cdots.
\end{equation}
\end{corollary}

\begin{proof}
By \eqref{eq:finite-trans-asymp} with \(K=0\), the difference from \(1\) is
\(\cO(q_1^n)\), where \(0<q_1<1\).  Exponential decay dominates every fixed
inverse power of \(n\).
\end{proof}

\begin{insightbox}
A pole of order \(s\) contributes a polynomial of degree \(s-1\) in \(n\)
times a geometric term.  Here every pole is simple, so \(s=1\): the polynomial
has degree zero.  This is the analytic reason that the factorial-normalized
Fubini transseries has one coefficient per pole and no \(1/n\) descendants.
The nonzero inverse-power coefficients that appear later come solely from
choosing to expand the exact factor \(n!\).
\end{insightbox}

\subsection{Coefficient dictionary in the natural transseries scale}

If the transmonomial of sector \(k\) is taken to be
\(
 \mathfrak m_k(n):=\e^{-(n+1)\Lambda_k}=\omega_k^{n+1}
\), then \eqref{eq:complex-transseries} reads
\begin{equation}\label{eq:natural-coeff}
 \frac{2\alpha^{n+1}}{n!}\Fub_n
 =\sum_{k\in\Z}\mathfrak m_k(n).
\end{equation}
Thus the full coefficient array is particularly simple:
\begin{equation}\label{eq:natural-array}
 c_{k,0}=1\quad(k\in\Z),
 \qquad
 c_{k,j}=0\quad(k\in\Z,\ j\ge1).
\end{equation}
This should not be mistaken for a lack of structure.  All geometry has moved
into the transmonomials themselves, through their exact complex actions
\(\Lambda_k\).

\section{Certified truncation and remainder bounds}

\subsection{The exact remainder after finitely many pole pairs}

Define the \(K\)-pair approximation
\begin{equation}\label{eq:SK}
 S_K(n):=\frac{n!}{2}\sum_{k=-K}^{K}\rho_k^{-n-1}
 =\frac{n!}{2\alpha^{n+1}}
 \left[1+2\sum_{k=1}^{K}q_k^{n+1}
 \cos\bigl((n+1)\theta_k\bigr)\right].
\end{equation}

\begin{theorem}[Hurwitz-zeta tail bound]\label{thm:tail}
For \(n\ge1\) and \(K\ge0\),
\begin{align}
 \abs{\Fub_n-S_K(n)}
 &\le n!\sum_{k=K+1}^{\infty}r_k^{-n-1}
 \label{eq:tail-exact-bound}\\
 &\le \frac{n!}{(2\pi)^{n+1}}
       \zeta(n+1,K+1),
 \label{eq:tail-hurwitz}
\end{align}
where \(\zeta(s,a)=\sum_{j=0}^{\infty}(j+a)^{-s}\) is the Hurwitz zeta
function.  In elementary terms,
\begin{equation}\label{eq:tail-elementary}
 \abs{\Fub_n-S_K(n)}
 \le \frac{n!}{(2\pi)^{n+1}}
 \left((K+1)^{-n-1}+\frac{(K+1)^{-n}}{n}\right).
\end{equation}
Relative to the dominant term
\(
 A_n=n!/(2\alpha^{n+1})
\),
\begin{equation}\label{eq:relative-tail}
 \frac{\abs{\Fub_n-S_K(n)}}{A_n}
 \le 2\left(\frac{\alpha}{2\pi}\right)^{n+1}
       \zeta(n+1,K+1).
\end{equation}
\end{theorem}

\begin{proof}
Pairing the omitted indices \(k\) and \(-k\) in
\eqref{eq:exact-poles} and applying the triangle inequality gives
\eqref{eq:tail-exact-bound}.  Since
\(
 r_k=\sqrt{\alpha^2+4\pi^2k^2}\ge2\pi k
\), equation \eqref{eq:tail-hurwitz} follows.  The integral test gives
\[
 \sum_{k=K+1}^{\infty}k^{-n-1}
 \le (K+1)^{-n-1}+\int_{K+1}^{\infty}x^{-n-1}\dd x,
\]
which is \eqref{eq:tail-elementary}.  Dividing
\eqref{eq:tail-hurwitz} by \(A_n\) gives \eqref{eq:relative-tail}.
\end{proof}

Taking \(K=0\) gives the explicit leading-term estimate
\begin{equation}\label{eq:leading-bound}
 \abs{\frac{\Fub_n}{A_n}-1}
 \le2\left(\frac{\alpha}{2\pi}\right)^{n+1}\zeta(n+1).
\end{equation}
This proves the classical equivalent \eqref{eq:leading} with a fully explicit,
exponentially small relative error.

\subsection{Relative accuracy versus absolute accuracy}

Because \(\Fub_n\) grows factorially, an approximation can have a spectacularly
small relative error and still miss the exact integer by a huge amount.  The
numerical values in \cref{tab:errors} illustrate this distinction.  Here
\(S_0\) uses only the real pole \(\rho_0\), while \(S_1\) also uses the first
conjugate pair.

\begin{table}[ht]
\centering
\caption{Absolute errors of pole truncations and relative error of the leading
term.  Values are rounded for display.}
\label{tab:errors}
\begin{tabular}{@{}rlll@{}}
\toprule
\(n\)&\(\abs{S_0(n)-\Fub_n}\)&\(\abs{S_1(n)-\Fub_n}\)&
\(\abs{S_0(n)/\Fub_n-1}\)\\
\midrule
5  &\(1.52\times10^{-3}\)&\(3.18\times10^{-5}\)&\(2.81\times10^{-6}\)\\
10 &\(5.27\times10^{-3}\)&\(1.66\times10^{-6}\)&\(5.16\times10^{-11}\)\\
16 &\(4.87\times10^{-1}\)&\(3.38\times10^{-6}\)&\(9.16\times10^{-17}\)\\
17 &\(5.42\times10^{-1}\)&\(3.10\times10^{-6}\)&\(4.16\times10^{-18}\)\\
20 &\(2.75\times10^{1}\)&\(1.78\times10^{-5}\)&\(1.03\times10^{-20}\)\\
30 &\(1.04\times10^{7}\)&\(2.11\times10^{-2}\)&\(9.10\times10^{-31}\)\\
50 &\(2.76\times10^{23}\)&\(7.98\times10^{7}\)&\(1.38\times10^{-49}\)\\
\bottomrule
\end{tabular}
\end{table}

The leading term rounds correctly through \(n=16\), but at \(n=17\) its
absolute error already exceeds \(1/2\).  This does not contradict asymptotic
equivalence: the relative error is then approximately \(4.16\times10^{-18}\).

\subsection{Certified nearest-integer recovery}

\begin{corollary}[A rigorous rounding certificate]\label{cor:rounding}
If
\begin{equation}\label{eq:rounding-condition}
 \frac{n!}{(2\pi)^{n+1}}
 \left((K+1)^{-n-1}+\frac{(K+1)^{-n}}{n}\right)<\frac12,
\end{equation}
then \(\Fub_n\) is the unique nearest integer to \(S_K(n)\).
\end{corollary}

\begin{proof}
The left side of \eqref{eq:rounding-condition} is an upper bound for
\(\abs{\Fub_n-S_K(n)}\) by \eqref{eq:tail-elementary}.  Since \(\Fub_n\) is an
integer, distance less than \(1/2\) identifies it uniquely by rounding.
\end{proof}

\begin{proposition}[A sufficient linear pole budget]\label{prop:linear-budget}
Fix \(\delta>0\).  If \(K=K_n\) is chosen so that
\begin{equation}\label{eq:linear-budget}
 K_n+1\ge \frac{(1+\delta)n}{2\pi\e},
\end{equation}
then the certified bound in \eqref{eq:tail-elementary} is
\begin{equation}\label{eq:linear-budget-error}
 \cO_\delta\!\left(n^{-1/2}(1+\delta)^{-n}\right).
\end{equation}
In particular, \eqref{eq:rounding-condition} holds for all sufficiently large
\(n\).
\end{proposition}

\begin{proof}
Use the elementary upper form of Stirling's estimate
\(
 n!\le C\sqrt n(n/\e)^n
\) with an absolute constant \(C\).  Put \(L=K_n+1\).  The two terms in
\eqref{eq:tail-elementary} are bounded by constant multiples of
\[
 \frac{\sqrt n}{L}
 \left(\frac{n}{2\pi\e L}\right)^n
 \quad\text{and}\quad
 \frac1{\sqrt n}
 \left(\frac{n}{2\pi\e L}\right)^n,
\]
respectively.  Condition \eqref{eq:linear-budget} makes the parenthesized ratio
at most \((1+\delta)^{-1}\).  It also gives
\(L\ge(1+\delta)n/(2\pi\e)\), hence
\(\sqrt n/L=\cO_\delta(n^{-1/2})\).  This proves
\eqref{eq:linear-budget-error}.
\end{proof}

Thus only about \(n/(2\pi\e)\approx0.05855n\) conjugate pole pairs are needed
for a simple asymptotic rounding certificate.  This is a sufficient asymptotic
budget, not a claim of pointwise minimality: oscillatory cancellation often
allows fewer pairs.

\section{The universal Stirling descendant tower}

\subsection{From the exact factorial prefactor to an elementary scale}

The natural formula \eqref{eq:exact-poles} keeps \(n!\) exact.  To rewrite it in
the more conventional elementary scale, define
\begin{equation}\label{eq:Pn}
 P_n:=\frac{\sqrt{2\pi n}}{2\alpha}
      \left(\frac{n}{\e\alpha}\right)^n.
\end{equation}
Stirling's expansion \cite[\S5.11]{DLMF} has the form
\begin{equation}\label{eq:stirling-gamma}
 n!\sim\sqrt{2\pi n}\left(\frac n\e\right)^n
 \sum_{j=0}^{\infty}\frac{\gamma_j}{n^j}.
\end{equation}
The coefficients \(\gamma_j\) are universal: they know nothing about Fubini
numbers.  We now derive them in a form suited to coefficient calculus.

\subsection{Bernoulli exponent and Bell-polynomial formula}

Use the Bernoulli convention
\begin{equation}\label{eq:bernoulli}
 \frac{t}{\e^t-1}=\sum_{m=0}^{\infty}B_m\frac{t^m}{m!},
 \qquad B_1=-\frac12.
\end{equation}
The logarithmic Stirling expansion is
\begin{equation}\label{eq:log-stirling}
 \log n!
 \sim\left(n+\frac12\right)\log n-n+\frac12\log(2\pi)
 +\sum_{r=1}^{\infty}\frac{a_r}{n^r},
\end{equation}
where
\begin{equation}\label{eq:ar}
 a_r:=\frac{B_{r+1}}{r(r+1)}.
\end{equation}
Thus \(a_1=1/12\), \(a_2=0\), \(a_3=-1/360\), and all even
\(a_r\) with \(r\ge2\) vanish.

Let \(Y_j(x_1,\ldots,x_j)\) denote the complete exponential Bell polynomial,
defined by
\begin{equation}\label{eq:bell-def}
 \exp\left(\sum_{r\ge1}x_r\frac{t^r}{r!}\right)
 =\sum_{j\ge0}Y_j(x_1,\ldots,x_j)\frac{t^j}{j!}.
\end{equation}

\begin{theorem}[General Stirling coefficient formulas]\label{thm:stirling-coeff}
The coefficients in \eqref{eq:stirling-gamma} are determined by
\begin{equation}\label{eq:gamma-gen}
 \sum_{j\ge0}\gamma_jt^j
 =\exp\left(\sum_{r\ge1}a_rt^r\right).
\end{equation}
Equivalently, for every \(j\ge0\),
\begin{align}
 \gamma_j
 &=\frac1{j!}Y_j(1!a_1,2!a_2,\ldots,j!a_j),
 \label{eq:gamma-bell}\\
 &=\sum_{\substack{m_1,\ldots,m_j\ge0\\
                    m_1+2m_2+\cdots+jm_j=j}}
   \prod_{r=1}^{j}\frac{a_r^{m_r}}{m_r!},
 \label{eq:gamma-partition}\\
 \gamma_0&=1,\qquad
 \gamma_j=\frac1j\sum_{r=1}^{j}
          \frac{B_{r+1}}{r+1}\gamma_{j-r}
          \quad(j\ge1).
 \label{eq:gamma-rec}
\end{align}
\end{theorem}

\begin{proof}
Exponentiating the correction series in \eqref{eq:log-stirling} gives
\eqref{eq:gamma-gen}.  Substitution \(x_r=r!a_r\) in
\eqref{eq:bell-def} gives \eqref{eq:gamma-bell}.  Expanding the exponential as
a product
\(
 \prod_{r\ge1}\exp(a_rt^r)
\)
and extracting \(t^j\) gives \eqref{eq:gamma-partition}.  Finally, logarithmic
differentiation of \eqref{eq:gamma-gen} gives
\[
 \frac{\sum_{j\ge1}j\gamma_jt^{j-1}}
      {\sum_{j\ge0}\gamma_jt^j}
 =\sum_{r\ge1}ra_rt^{r-1}.
\]
Comparing coefficients of \(t^{j-1}\) and using
\(
 ra_r=B_{r+1}/(r+1)
\) proves \eqref{eq:gamma-rec}.
\end{proof}

The first coefficients are listed in \cref{tab:stirling-coeffs}.

\begin{table}[ht]
\centering
\caption{Stirling coefficients \(\gamma_j\) in \eqref{eq:stirling-gamma}.}
\label{tab:stirling-coeffs}
{\small
\renewcommand{\arraystretch}{2.65}
\begin{tabular}{@{}r c @{\hspace{3.2em}} r c@{}}
\toprule
\(j\)&\(\displaystyle \gamma_j\)&\(j\)&\(\displaystyle \gamma_j\)\\
\midrule
0&\(\displaystyle 1\)
 &6&\(\displaystyle \frac{5246819}{75246796800}\)\\
1&\(\displaystyle \frac1{12}\)
 &7&\(\displaystyle -\frac{534703531}{902961561600}\)\\
2&\(\displaystyle \frac1{288}\)
 &8&\(\displaystyle -\frac{4483131259}{86684309913600}\)\\
3&\(\displaystyle -\frac{139}{51840}\)
 &9&\(\displaystyle \frac{432261921612371}{514904800886784000}\)\\
4&\(\displaystyle -\frac{571}{2488320}\)
 &10&\(\displaystyle \frac{6232523202521089}{86504006548979712000}\)\\
5&\(\displaystyle \frac{163879}{209018880}\)
 &&\\
\bottomrule
\end{tabular}}
\end{table}

\subsection{The full two-level coefficient theorem}

\begin{theorem}[All pole-sector and algebraic coefficients]\label{thm:double-transseries}
Let \(P_n\) be given by \eqref{eq:Pn}, and let
\(
 \Lambda_k=\Log(\rho_k/\alpha)
\).
Then the Fubini numbers have the formal two-level transseries
\begin{equation}\label{eq:double-transseries}
 \boxed{\displaystyle
 \Fub_n\sim P_n
 \sum_{k\in\Z}\e^{-n\Lambda_k}
 \sum_{j=0}^{\infty}\frac{C_{k,j}}{n^j}}
\end{equation}
with the explicit general coefficient
\begin{equation}\label{eq:Ckj}
 \boxed{\displaystyle
 C_{k,j}=\e^{-\Lambda_k}\gamma_j
 =\frac{\alpha}{\alpha+2\pi\ii k}\,
   \frac1{j!}Y_j(1!a_1,\ldots,j!a_j)}.
\end{equation}
Equivalently,
\begin{equation}\label{eq:Ckj-partition}
 C_{k,j}=\frac{\alpha}{\alpha+2\pi\ii k}
 \sum_{\substack{m_1,\ldots,m_j\ge0\\
                  \sum_{r=1}^jrm_r=j}}
 \prod_{r=1}^{j}\frac1{m_r!}
 \left(\frac{B_{r+1}}{r(r+1)}\right)^{m_r}.
\end{equation}
\end{theorem}

\begin{proof}
Write \eqref{eq:exact-poles} as
\[
 \Fub_n=\frac{n!}{2\alpha^{n+1}}
 \sum_{k\in\Z}\e^{-(n+1)\Lambda_k}.
\]
Insert \eqref{eq:stirling-gamma}, and observe that
\[
 \frac{\sqrt{2\pi n}(n/\e)^n}{2\alpha^{n+1}}=P_n,
 \qquad
 \e^{-(n+1)\Lambda_k}
 =\e^{-n\Lambda_k}\e^{-\Lambda_k}.
\]
The coefficient of \(\e^{-n\Lambda_k}n^{-j}\) is therefore
\(\e^{-\Lambda_k}\gamma_j\), which is \eqref{eq:Ckj}.  Substituting
\eqref{eq:gamma-partition} gives \eqref{eq:Ckj-partition}.
\end{proof}

Pairing conjugates gives the real form
\begin{equation}\label{eq:double-real}
 \Fub_n\sim P_n
 \left(\sum_{j\ge0}\frac{\gamma_j}{n^j}\right)
 \left[1+2\sum_{k\ge1}q_k^{n+1}
 \cos\bigl((n+1)\theta_k\bigr)\right].
\end{equation}
The bracketed factor is exact; only the Stirling factor is asymptotic.

\subsection{A precise double-truncation statement}

For integers \(J\ge1\) and \(K\ge0\), define
\begin{equation}\label{eq:TJK}
 T_{J,K}(n):=P_n
 \left(\sum_{j=0}^{J-1}\frac{\gamma_j}{n^j}\right)
 \sum_{k=-K}^{K}\e^{-(n+1)\Lambda_k}.
\end{equation}
For fixed \(J,K\), the standard remainder statement for Stirling's expansion
and \cref{thm:tail} give
\begin{align}
 \Fub_n-T_{J,K}(n)
 &={}\cO_{J,K}\!\left(P_n n^{-J}\right)
 \nonumber\\
 &\quad+\cO\!\left(
       \frac{n!}{(2\pi)^{n+1}}\zeta(n+1,K+1)
       \right).
 \label{eq:double-remainder}
\end{align}
The first term is the algebraic truncation error of the exact factorial; the
second is the omitted Fubini pole tail.  The two truncations are independent and
should be chosen according to the desired task.  For relative asymptotics, a
small fixed \(K\) is usually enough.  For exact integer reconstruction, \(K\)
must grow with \(n\) as discussed in \cref{prop:linear-budget}.

\section{A multi-index coefficient calculus for nonlinear observables}

The linear sequence \(\Fub_n\) itself contains one copy of each pole sector.
Products, powers, logarithms, reciprocals, and ratios generate sums of pole
actions.  This section gives their coefficients in closed form.

\subsection{The normalized exponential correction}

Put
\begin{equation}\label{eq:HnEn}
 H_n:=\frac{2\alpha^{n+1}}{n!}\Fub_n
 =1+E_n,
 \qquad
 E_n:=\sum_{k\ne0}\e^{-(n+1)\Lambda_k}.
\end{equation}
By \cref{cor:flat}, \(E_n\to0\) exponentially.
Let \(\mathcal I\) be the set of finitely supported multi-indices
\(
 \bm m=(m_k)_{k\in\Z\setminus\{0\}}
\)
with \(m_k\in\N\).  Define
\begin{equation}\label{eq:multi-def}
 \abs{\bm m}:=\sum_{k\ne0}m_k,
 \qquad
 \Lambda_{\bm m}:=\sum_{k\ne0}m_k\Lambda_k,
 \qquad
 \bm m!:=\prod_{k\ne0}m_k!.
\end{equation}
All sums and products are finite for each \(\bm m\).

\begin{theorem}[Analytic-function coefficient rule]\label{thm:analytic-phi}
Let \(\Phi\) be analytic in a neighborhood of \(1\).  Formally, and
analytically for all sufficiently large \(n\),
\begin{equation}\label{eq:phi-multi}
 \Phi(H_n)
 =\sum_{\bm m\in\mathcal I}
   \frac{\Phi^{(\abs{\bm m})}(1)}{\bm m!}
   \e^{-(n+1)\Lambda_{\bm m}}.
\end{equation}
The zero multi-index contributes \(\Phi(1)\).
\end{theorem}

\begin{proof}
Taylor-expand \(\Phi(1+E_n)\):
\[
 \Phi(1+E_n)=\sum_{\ell\ge0}
 \frac{\Phi^{(\ell)}(1)}{\ell!}E_n^\ell.
\]
The multinomial theorem gives
\[
 E_n^\ell
 =\sum_{\substack{\bm m\in\mathcal I\\\abs{\bm m}=\ell}}
   \frac{\ell!}{\bm m!}
   \e^{-(n+1)\Lambda_{\bm m}}.
\]
Substitution proves \eqref{eq:phi-multi}.  Since \(E_n\to0\), the Taylor series
converges analytically for sufficiently large \(n\).
\end{proof}

This theorem is the transseries analogue of the ordinary multinomial
coefficient calculus.  It also displays a useful distinction: the primitive
Fubini sequence has only single-pole sectors, while nonlinear observables have
all additive combinations \(\Lambda_{\bm m}\).

\subsection{Arbitrary complex powers}

For \(\Phi(u)=u^\beta\), where a branch is fixed near \(u=1\),
\(
 \Phi^{(\ell)}(1)=\falling{\beta}{\ell}
\).  Hence
\begin{equation}\label{eq:Hbeta}
 H_n^\beta
 =\sum_{\bm m\in\mathcal I}
   \frac{\falling{\beta}{\abs{\bm m}}}{\bm m!}
   \e^{-(n+1)\Lambda_{\bm m}}.
\end{equation}
For \(\beta=1\), the falling factorial vanishes for
\(\abs{\bm m}\ge2\); this recovers the fact that \(H_n\) contains only the
vacuum and one-pole sectors.

The algebraic Stirling factor also has an all-orders power calculus.  Define
\begin{equation}\label{eq:gamma-beta-gen}
 \sum_{j\ge0}\gamma_j^{(\beta)}t^j
 :=\exp\left(\beta\sum_{r\ge1}a_rt^r\right).
\end{equation}
Then
\begin{align}
 \gamma_j^{(\beta)}
 &=\frac1{j!}Y_j(\beta1!a_1,\ldots,\beta j!a_j),
 \label{eq:gamma-beta-bell}\\
 &=\sum_{\substack{m_1,\ldots,m_j\ge0\\\sum rm_r=j}}
   \beta^{m_1+\cdots+m_j}
   \prod_{r=1}^{j}\frac{a_r^{m_r}}{m_r!}.
 \label{eq:gamma-beta-partition}
\end{align}
Combining \eqref{eq:Hbeta} and \eqref{eq:gamma-beta-gen} yields the complete
power transseries.

\begin{corollary}[All coefficients of \(\Fub_n^\beta\)]\label{cor:Fbeta}
With \(P_n\) as in \eqref{eq:Pn},
\begin{equation}\label{eq:Fbeta}
 \Fub_n^\beta\sim P_n^\beta
 \left(\sum_{j\ge0}\frac{\gamma_j^{(\beta)}}{n^j}\right)
 \sum_{\bm m\in\mathcal I}
 \frac{\falling{\beta}{\abs{\bm m}}}{\bm m!}
 \e^{-(n+1)\Lambda_{\bm m}}.
\end{equation}
Consequently, the coefficient at algebraic order \(j\) and composite action
\(\Lambda_{\bm m}\) is
\begin{equation}\label{eq:Fbeta-coeff}
 \boxed{\displaystyle
 C_{\bm m,j}^{(\beta)}
 =\gamma_j^{(\beta)}
  \frac{\falling{\beta}{\abs{\bm m}}}{\bm m!}}.
\end{equation}
If the transmonomial is written as \(\e^{-n\Lambda_{\bm m}}\) rather than
\(\e^{-(n+1)\Lambda_{\bm m}}\), the coefficient is multiplied by
\(\e^{-\Lambda_{\bm m}}\).
\end{corollary}

\subsection{The logarithm}

For \(\Phi(u)=\log u\),
\(
 \Phi^{(\ell)}(1)=(-1)^{\ell-1}(\ell-1)!
\) for \(\ell\ge1\).  Thus
\begin{equation}\label{eq:logH}
 \log H_n
 =\sum_{\substack{\bm m\in\mathcal I\\\bm m\ne\bm0}}
 (-1)^{\abs{\bm m}+1}
 \frac{(\abs{\bm m}-1)!}{\bm m!}
 \e^{-(n+1)\Lambda_{\bm m}}.
\end{equation}
Since the logarithm of the Stirling factor is exactly the Bernoulli exponent in
\eqref{eq:log-stirling}, we obtain
\begin{align}
 \log\Fub_n
 \sim{}&\log P_n+\sum_{r\ge1}\frac{a_r}{n^r}
 \nonumber\\
 &+\sum_{\substack{\bm m\in\mathcal I\\\bm m\ne\bm0}}
 (-1)^{\abs{\bm m}+1}
 \frac{(\abs{\bm m}-1)!}{\bm m!}
 \e^{-(n+1)\Lambda_{\bm m}}.
 \label{eq:logF-full}
\end{align}
Equations \eqref{eq:ar} and \eqref{eq:logF-full} give every algebraic and every
composite exponential coefficient of \(\log\Fub_n\).

\subsection{Reciprocals and ratios}

For the reciprocal, \(\Phi(u)=u^{-1}\), equation
\eqref{eq:phi-multi} gives
\begin{equation}\label{eq:reciprocal-H}
 \frac1{H_n}
 =\sum_{\bm m\in\mathcal I}
 (-1)^{\abs{\bm m}}
 \frac{\abs{\bm m}!}{\bm m!}
 \e^{-(n+1)\Lambda_{\bm m}}.
\end{equation}
An exact ratio identity follows directly from \eqref{eq:HnEn}:
\begin{equation}\label{eq:ratio-exact}
 \frac{n\Fub_{n-1}}{\Fub_n}
 =\alpha\frac{H_{n-1}}{H_n}.
\end{equation}
Since
\(
 H_{n-1}=1+\sum_{k\ne0}\omega_k^n
\), linearization gives
\begin{equation}\label{eq:ratio-linear}
 \frac{n\Fub_{n-1}}{\alpha\Fub_n}
 =1+\sum_{k\ne0}(1-\omega_k)\omega_k^n
 +\cO(q_1^{2n}).
\end{equation}
Keeping only the nearest conjugate pair gives
\begin{align}
 \frac{n\Fub_{n-1}}{\alpha\Fub_n}
 ={}&1+2q_1^n\Re\left((1-\omega_1)\e^{-\ii n\theta_1}\right)
 +\cO(q_2^n).
 \label{eq:ratio-first}
\end{align}
In particular,
\(
 n\Fub_{n-1}/\Fub_n\to\log2
\), with an explicitly oscillatory exponential approach.

\section{Borel transform and resurgent large-order structure}

\subsection{The divergent ordinary series and its exact Borel transform}

The ordinary generating series
\begin{equation}\label{eq:formal-ogf}
 \widehat{\mathcal F}(x):=\sum_{n=0}^{\infty}\Fub_nx^{n+1}
\end{equation}
has zero radius of convergence because \(\Fub_n\) grows factorially.  It is,
however, Gevrey of order one.  With the convention
\begin{equation}\label{eq:borel-def}
 \Borel\left(\sum_{n\ge0}a_nx^{n+1}\right)(\xi)
 :=\sum_{n\ge0}a_n\frac{\xi^n}{n!},
\end{equation}
we obtain immediately from \eqref{eq:basic-egf}:

\begin{theorem}[Exact Borel transform]\label{thm:borel}
\begin{equation}\label{eq:borel-exact}
 \boxed{\displaystyle
 \Borel\widehat{\mathcal F}(\xi)=\frac1{2-\e^\xi}}.
\end{equation}
It is meromorphic on \(\C\), with simple singularities
\(\xi=\rho_k\) and residue \(-1/2\) at every one.
\end{theorem}

The singularity dictionary is therefore exact:
\begin{center}
\begin{tabular}{@{}lll@{}}
\toprule
Borel datum&value&large-order contribution\\
\midrule
singularity&\(\rho_k=\alpha+2\pi\ii k\)&sector \(k\)\\
residue&\(-1/2\)&\(+n!/(2\rho_k^{n+1})\)\\
relative action&\(\Lambda_k=\Log(\rho_k/\alpha)\)&\(\e^{-(n+1)\Lambda_k}\)\\
\bottomrule
\end{tabular}
\end{center}
Summing the large-order contribution over all Borel singularities reproduces
\eqref{eq:exact-poles} exactly.  In this example, the usual resurgent
``large-order relation'' closes without an uncomputed remainder.

\subsection{Lateral Borel--Laplace sums and Stokes jumps}

For a direction \(\vartheta\) not containing a pole, define formally
\begin{equation}\label{eq:laplace}
 \mathcal S_\vartheta\widehat{\mathcal F}(x)
 :=\int_{0}^{\e^{\ii\vartheta}\infty}
 \e^{-\xi/x}\frac{\dd\xi}{2-\e^\xi},
\end{equation}
whenever the integral converges.  Expanding the Borel transform at the origin
and integrating term by term recovers \eqref{eq:formal-ogf} asymptotically.
The positive real direction is a Stokes direction because it contains
\(\rho_0=\alpha\).  With \(\mathcal S_{0^+}\) and
\(\mathcal S_{0^-}\) denoting paths just above and below the positive axis,
our orientation convention gives
\begin{equation}\label{eq:stokes-jump}
 \mathcal S_{0^+}\widehat{\mathcal F}(x)
 -\mathcal S_{0^-}\widehat{\mathcal F}(x)
 =\pi\ii\e^{-\alpha/x}.
\end{equation}
Indeed, the jump is \(-2\pi\ii\) times the residue of
\(
 \e^{-\xi/x}(2-\e^\xi)^{-1}
\) at \(\xi=\alpha\), and that residue is
\(-\tfrac12\e^{-\alpha/x}\).  Analogous jumps occur on rays through the other
\(\rho_k\).  Standard background on Borel summation and resurgence can be
found in Costin~\cite{Costin} and Sauzin~\cite{Sauzin}.

\begin{remark}[Two different transseries variables]
The large-order transseries studied in the preceding sections uses
\(n\to\infty\) and actions \(\Lambda_k\).  The Borel--Laplace problem uses
\(x\to0\) and actions \(\rho_k/x\).  They encode the same pole lattice in two
dual ways: \(\rho_k^{-n-1}\) controls Taylor-coefficient growth, while
\(\e^{-\rho_k/x}\) controls Stokes phenomena of the Borel sum.
\end{remark}

\section{Fubini polynomials and a one-parameter pole lattice}

\subsection{Definition and geometric-moment representation}

Define the Fubini polynomials
\begin{equation}\label{eq:fubini-poly-def}
 \Fub_n(y):=\sum_{j=0}^n j!\Stirling{n}{j}y^j.
\end{equation}
Their exponential generating function is
\begin{equation}\label{eq:fubini-poly-egf}
 \sum_{n\ge0}\Fub_n(y)\frac{z^n}{n!}
 =\frac1{1-y(\e^z-1)}
 =\frac1{1+y-y\e^z}.
\end{equation}
For \(y>0\), set
\begin{equation}\label{eq:alpha-y}
 \alpha_y:=\log\left(1+\frac1y\right),
 \qquad
 \rho_k(y):=\alpha_y+2\pi\ii k.
\end{equation}
Expanding \eqref{eq:fubini-poly-egf} geometrically gives
\begin{equation}\label{eq:fubini-poly-moment}
 \Fub_n(y)=\frac1{1+y}
 \sum_{m=0}^{\infty}
 \left(\frac{y}{1+y}\right)^m m^n.
\end{equation}

\begin{theorem}[Complete transseries of the Fubini polynomials]\label{thm:fubini-poly}
For \(y>0\) and \(n\ge1\),
\begin{equation}\label{eq:fubini-poly-poles}
 \boxed{\displaystyle
 \Fub_n(y)=\frac{n!}{1+y}
 \sum_{k\in\Z}\rho_k(y)^{-n-1}}.
\end{equation}
Consequently,
\begin{align}
 \Fub_n(y)
 ={}&\frac{n!}{(1+y)\alpha_y^{n+1}}
 \Biggl[1+2\sum_{k\ge1}
 \left(\frac{\alpha_y}
 {\sqrt{\alpha_y^2+4\pi^2k^2}}\right)^{n+1}
 \nonumber\\[-1mm]
 &\hspace{31mm}\times
 \cos\left((n+1)\arctan\frac{2\pi k}{\alpha_y}\right)
 \Biggr].
 \label{eq:fubini-poly-real}
\end{align}
All coefficient and remainder formulas of the preceding sections remain valid
after replacing \(\alpha\) by \(\alpha_y\) and \(1/2\) by \(1/(1+y)\).
\end{theorem}

\begin{proof}
Equation \eqref{eq:fubini-poly-moment} is the moment sum with decay parameter
\(\alpha_y\), because
\(
 y/(1+y)=\e^{-\alpha_y}
\).  Applying the proof of \cref{thm:exact-poles} gives
\eqref{eq:fubini-poly-poles}; pairing conjugates gives
\eqref{eq:fubini-poly-real}.
\end{proof}

The ordinary Fubini numbers are recovered at \(y=1\), for which
\(\alpha_y=\log2\) and \(1/(1+y)=1/2\).

\section{Higher-order poles and N\"orlund coefficient calculus}

The simplicity of the Fubini transseries is caused by simple poles.  It is
useful to see the general pattern when those poles are raised to order \(r\).
This also supplies a nontrivial family of all-orders pole coefficients.

\subsection{Convolution powers}

For an integer \(r\ge1\), define
\begin{equation}\label{eq:conv-power}
 A_n^{(r)}:=n!\coeff{z}{n}(2-\e^z)^{-r}.
\end{equation}
Combinatorially, \(A_n^{(r)}\) is the \(r\)-fold binomial convolution of the
Fubini sequence, because products of exponential generating functions encode
binomial convolution.

Define the generalized Bernoulli, or N\"orlund, numbers by
\begin{equation}\label{eq:norlund}
 \left(\frac{t}{\e^t-1}\right)^r
 =\sum_{j=0}^{\infty}B_j^{(r)}\frac{t^j}{j!}.
\end{equation}
For example,
\(
 B_0^{(r)}=1
\),
\(
 B_1^{(r)}=-r/2
\), and
\(
 B_2^{(r)}=r(3r-1)/12
\).

\begin{theorem}[General coefficient formula for order-\(r\) pole sectors]
\label{thm:higher-poles}
For \(r\ge1\) and \(n\ge1\),
\begin{equation}\label{eq:higher-pole-formula}
 \boxed{\displaystyle
 A_n^{(r)}=
 \frac{n!}{2^r}\sum_{k\in\Z}\sum_{j=0}^{r-1}
 \frac{(-1)^jB_j^{(r)}}{j!}
 \binom{n+r-j-1}{r-j-1}
 \rho_k^{-n-r+j}.}
\end{equation}
Thus each pole sector is a polynomial in \(n\) of degree \(r-1\), multiplied
by the geometric factor \(\rho_k^{-n}\).  Formula \eqref{eq:exact-poles} is the
case \(r=1\).
\end{theorem}

\begin{proof}
We first justify that the coefficient is determined exactly by the principal
parts at the full pole lattice.  Apply the residue theorem to
\[
 \frac{(2-\e^z)^{-r}}{z^{n+1}}
\]
on rectangles whose vertical sides are \(\Re z=\pm R\) and whose horizontal
sides are \(\Im z=\pm 2\pi(N+\tfrac12)\), with \(R,N\to\infty\) in a
comparable way.  On the left side the denominator \(2-\e^z\) stays bounded
away from zero; on the right side its reciprocal decays exponentially; and on
the horizontal sides \(\e^z=-\e^{\Re z}\), so again the denominator stays
away from zero.  Since \(n\ge1\), the four boundary integrals therefore tend
to zero.  Consequently, the residue at the origin, which is
\(\coeff{z}{n}(2-\e^z)^{-r}\), equals the negative of the sum of the residues
at all \(\rho_k\).

Near a pole \(\rho_k\), put \(t=z-\rho_k\).  Since
\(
 \e^{\rho_k}=2
\),
\begin{align}
 (2-\e^z)^{-r}
 &=(-2)^{-r}(\e^t-1)^{-r}
 \nonumber\\
 &=(-2)^{-r}t^{-r}
   \left(\frac{t}{\e^t-1}\right)^r
 \nonumber\\
 &=(-2)^{-r}
   \sum_{j\ge0}\frac{B_j^{(r)}}{j!}t^{j-r}.
 \label{eq:local-r}
\end{align}
Only \(j=0,1,\ldots,r-1\) belong to the principal part.  If
\(m=r-j\), then
\begin{equation}\label{eq:coeff-pole-m}
 \coeff{z}{n}(z-\rho_k)^{-m}
 =(-1)^m\binom{n+m-1}{m-1}\rho_k^{-n-m}.
\end{equation}
This coefficient is precisely the negative of the residue contribution of
\((z-\rho_k)^{-m}z^{-n-1}\), in agreement with the preceding contour
identity.  Multiplying \eqref{eq:coeff-pole-m} by the coefficient of
\(t^{-m}\) in \eqref{eq:local-r} gives the inner summand of
\eqref{eq:higher-pole-formula}; the signs simplify to \((-1)^j\).  Finally,
every term is \(\cO(|k|^{-n-1})\) or smaller, so the pole sum is absolutely
convergent for \(n\ge1\), completing the proof.
\end{proof}

For \(r=2\), one has \(B_0^{(2)}=1\) and \(B_1^{(2)}=-1\), so
\begin{equation}\label{eq:r2}
 A_n^{(2)}=\frac{n!}{4}\sum_{k\in\Z}
 \left[(n+1)\rho_k^{-n-2}+\rho_k^{-n-1}\right].
\end{equation}
The degree-one polynomial in \(n\) is exactly what a double pole predicts.
This comparison isolates the special feature of the original Fubini sequence:
its degree-zero sector polynomial is not an approximation but the whole local
contribution.

\section{Numerical verification and reproducible computation}

\subsection{Exact arithmetic checks}

For \(n\ge1\), the three exact representations
\begin{equation}\label{eq:three-exact}
 \Fub_n
 =\sum_{j=0}^n j!\Stirling{n}{j}
 =\frac12\sum_{m\ge0}\frac{m^n}{2^m}
 =\frac{n!}{2}\sum_{k\in\Z}\rho_k^{-n-1}
\end{equation}
serve different computational purposes.

\begin{itemize}
\item The finite Stirling sum is ideal for exact integer computation.
\item The geometric moment sum is positive and conceptually simple, but its
      efficient truncation depends on locating the maximum near \(m=n/\alpha\).
\item The pole sum converges extraordinarily fast in relative terms as \(n\)
      grows and comes with the certified bound \eqref{eq:tail-hurwitz}.
\end{itemize}

The following Wolfram Language fragment implements the finite definition, the
pole approximation, and the rigorous elementary tail bound.

\begin{lstlisting}[language=Mathematica]
ClearAll[fubini, rho, poleApproximation, poleTailBound];

fubini[n_Integer?NonNegative] :=
  Sum[k! StirlingS2[n, k], {k, 0, n}];

rho[k_Integer] := Log[2] + 2 Pi I k;

poleApproximation[n_Integer?Positive, K_Integer?NonNegative] :=
  n!/2 Sum[rho[k]^(-n - 1), {k, -K, K}];

poleTailBound[n_Integer?Positive, K_Integer?NonNegative] :=
  n!/(2 Pi)^(n + 1) ((K + 1)^(-n - 1) +
    (K + 1)^(-n)/n);
\end{lstlisting}

A certified exact-recovery routine increases \(K\) until
\(
 \texttt{poleTailBound[n,K]}<1/2
\), evaluates the symmetric pole sum with enough guard digits, and rounds its
real part.  The imaginary part should vanish up to numerical error because the
summands occur in conjugate pairs.

\subsection{Stable real evaluation}

For numerical work it is preferable to use the explicitly real formula
\eqref{eq:SK} rather than summing complex powers independently.  A stable
implementation is
\begin{lstlisting}[language=Mathematica]
ClearAll[realPoleApproximation];

realPoleApproximation[n_Integer?Positive,
                      K_Integer?NonNegative] := Module[
  {a = Log[2]},
  n!/(2 a^(n + 1)) (1 +
    2 Sum[
      (a/Sqrt[a^2 + 4 Pi^2 k^2])^(n + 1) *
      Cos[(n + 1) ArcTan[2 Pi k/a]],
      {k, 1, K}])
  ];
\end{lstlisting}
The scale separation in \(q_k^{n+1}\) means that high working precision may be
necessary when the goal is exact rounding rather than relative approximation.
One should request enough digits to resolve the absolute tail target, not just
the relative size of \(\Fub_n\).

\subsection{Generation of Stirling descendants}

The recurrence \eqref{eq:gamma-rec} gives a direct exact algorithm:
\begin{lstlisting}[language=Mathematica]
ClearAll[stirlingGamma];
stirlingGamma[0] = 1;
stirlingGamma[j_Integer?Positive] := stirlingGamma[j] =
  1/j Sum[BernoulliB[r + 1]/(r + 1) *
          stirlingGamma[j - r], {r, 1, j}];
\end{lstlisting}
Alternatively, the complete Bell-polynomial formula
\eqref{eq:gamma-bell} can be used directly.  Both calculations are exact over
the rationals.

\section{Consequences, comparisons, and interpretation}

\subsection{A complete hierarchy of approximants}

The exact real transseries gives the nested sequence
\begin{align*}
 \Fub_n
 &\sim \frac{n!}{2\alpha^{n+1}},\\
 \Fub_n
 &\sim \frac{n!}{2\alpha^{n+1}}
 \left[1+2q_1^{n+1}\cos((n+1)\theta_1)\right],\\
 \Fub_n
 &\sim \frac{n!}{2\alpha^{n+1}}
 \left[1+2\sum_{k=1}^{2}q_k^{n+1}
 \cos((n+1)\theta_k)\right],
\end{align*}
and so on.  Unlike a divergent power series, this hierarchy is obtained by
truncating an absolutely convergent exact sum.  There is no least term and no
optimal finite truncation for relative accuracy: adding a pole pair always
reduces the rigorous absolute tail bound.  Optimality enters only through
computational cost or through a simultaneous truncation of the divergent
Stirling series.

\subsection{Contrast with ordinary Bell numbers}

The ordinary Bell numbers have entire exponential generating function
\(
 \exp(\e^z-1)
\).  Their asymptotics arise from moving saddle points and naturally involve
the Lambert \(W\)-function and nontrivial algebraic correction towers.  Fubini
numbers instead have a meromorphic generating function with a fixed, explicitly
known pole lattice.  The difference is decisive:

\begin{center}
\begin{tabularx}{\textwidth}{@{}>{\bfseries}l>{\raggedright\arraybackslash}X>{\raggedright\arraybackslash}X@{}}
\toprule
&Ordinary Bell numbers&Fubini numbers\\
\midrule
EGF&\(\exp(\e^z-1)\), entire&\((2-\e^z)^{-1}\), meromorphic\\
Mechanism&moving saddle point&fixed singularity lattice\\
Natural special function&Lambert \(W\)&elementary logarithms\\
Normalized corrections&nontrivial powers of \(1/n\)&pure exponentials, no power tower\\
Full sector data&saddle expansion&exact residues of all poles\\
\bottomrule
\end{tabularx}
\end{center}

This makes the Fubini sequence a particularly clean model for learning how
singularity lattices generate transseries.

\subsection{Completeness of the result}

The word ``complete'' can be justified concretely.

\begin{enumerate}
\item The Borel transform \((2-\e^\xi)^{-1}\) is known exactly on the whole
      plane and has no singularities other than \(\rho_k\).
\item Every singularity is included in \eqref{eq:exact-poles} with its exact
      residue contribution.
\item The pole sum is absolutely convergent for every \(n\ge1\), so no
      unspecified summation prescription is needed.
\item The conventional inverse-power descendants introduced by expanding
      \(n!\) have the explicit coefficient formulas
      \eqref{eq:gamma-bell}--\eqref{eq:gamma-rec}.
\item All additive composite sectors created by nonlinear operations have the
      multi-index coefficients in \eqref{eq:phi-multi},
      \eqref{eq:Fbeta-coeff}, and \eqref{eq:logH}.
\item Every finite pole truncation has the certified remainder
      \eqref{eq:tail-hurwitz}.
\end{enumerate}

There is therefore no hidden Fubini-specific algebraic series, no unlisted
exponential action, and no unquantified pole tail.

\section{Summary of principal formulae}

For convenient reference, the main formulas are collected here.

\begin{notationbox}
\[
 \alpha=\log2,
 \quad \rho_k=\alpha+2\pi\ii k,
 \quad \omega_k=\frac{\alpha}{\rho_k},
 \quad \Lambda_k=\Log\frac{\rho_k}{\alpha}.
\]
For \(k\ge1\),
\[
 q_k=\frac{\alpha}{\sqrt{\alpha^2+4\pi^2k^2}},\quad
 \lambda_k=-\log q_k,\quad
 \theta_k=\arctan\frac{2\pi k}{\alpha}.
\]
\end{notationbox}

\begin{align}
 \Fub_n
 &=\sum_{j=0}^n j!\Stirling{n}{j},
 \label{eq:summary1}\\
 \sum_{n\ge0}\Fub_n\frac{z^n}{n!}
 &=\frac1{2-\e^z},
 \label{eq:summary2}\\
 \Fub_n
 &=\frac12\sum_{m\ge0}\frac{m^n}{2^m},
 \label{eq:summary3}\\
 \Fub_n
 &=\frac{n!}{2}\sum_{k\in\Z}\rho_k^{-n-1},
 \qquad n\ge1,
 \label{eq:summary4}\\
 \frac{2\alpha^{n+1}}{n!}\Fub_n
 &=1+2\sum_{k\ge1}q_k^{n+1}
 \cos((n+1)\theta_k),
 \label{eq:summary5}\\
 \abs{\Fub_n-S_K(n)}
 &\le \frac{n!}{(2\pi)^{n+1}}\zeta(n+1,K+1),
 \label{eq:summary6}\\
 \Fub_n
 &\sim \frac{\sqrt{2\pi n}}{2\alpha}
 \left(\frac{n}{\e\alpha}\right)^n
 \sum_{k\in\Z}\e^{-n\Lambda_k}
 \sum_{j\ge0}\frac{C_{k,j}}{n^j},
 \label{eq:summary7}\\
 C_{k,j}
 &=\frac{\alpha}{\rho_k}\,
 \frac1{j!}Y_j(1!a_1,\ldots,j!a_j),
 \qquad
 a_r=\frac{B_{r+1}}{r(r+1)},
 \label{eq:summary8}\\
 \Phi(H_n)
 &=\sum_{\bm m\in\mathcal I}
 \frac{\Phi^{(\abs{\bm m})}(1)}{\bm m!}
 \e^{-(n+1)\Lambda_{\bm m}},
 \label{eq:summary9}\\
 \Borel\widehat{\mathcal F}(\xi)
 &=\frac1{2-\e^\xi}.
 \label{eq:summary10}
\end{align}

\appendix

\section{Derivation of the Bell-polynomial coefficient formulas}

For completeness, we spell out the coefficient mechanism behind
\cref{thm:stirling-coeff}.  If
\(
 A(t)=\sum_{r\ge1}a_rt^r
\), then
\[
 \e^{A(t)}
 =\prod_{r\ge1}\sum_{m_r\ge0}\frac{a_r^{m_r}t^{rm_r}}{m_r!}.
\]
The coefficient of \(t^j\) is obtained by imposing
\(
 \sum_{r=1}^jrm_r=j
\), which gives \eqref{eq:gamma-partition}.  This is a partition of the
integer \(j\) in frequency form: \(m_r\) records how many parts of size \(r\)
appear.  The complete exponential Bell polynomial packages the same sum after
the rescaling \(x_r=r!a_r\).

Because \(a_{2s}=0\) for \(s\ge1\), the partition sum may be reduced to odd
parts:
\begin{equation}\label{eq:gamma-odd}
 \gamma_j=
 \sum_{\substack{m_1,m_3,m_5,\ldots\ge0\\
                  \sum_{s\ge1}(2s-1)m_{2s-1}=j}}
 \prod_{s\ge1}\frac1{m_{2s-1}!}
 \left(\frac{B_{2s}}{2s(2s-1)}\right)^{m_{2s-1}}.
\end{equation}
Only finitely many indices can occur for fixed \(j\).  Formula
\eqref{eq:gamma-odd} is often the most efficient closed form for symbolic
manipulation.

For example,
\begin{align*}
 \gamma_3
 &=a_3+\frac{a_1^3}{3!}
 =-\frac1{360}+\frac1{6\cdot12^3}
 =-\frac{139}{51840},\\
 \gamma_4
 &=a_1a_3+\frac{a_1^4}{4!}
 =-\frac1{4320}+\frac1{24\cdot12^4}
 =-\frac{571}{2488320}.
\end{align*}
These calculations illustrate how Bernoulli data in the logarithm are converted
into ordinary Stirling coefficients by set-partition combinatorics.

\section{A coefficient-extraction proof of a pole contribution}

Suppose a meromorphic generating function has a pole principal part
\(
 A(z-\rho)^{-m}
\), where \(m\ge1\).  Then
\begin{align*}
 (z-\rho)^{-m}
 &=(-\rho)^{-m}\left(1-\frac z\rho\right)^{-m}\\
 &=(-1)^m\rho^{-m}
   \sum_{n\ge0}\binom{n+m-1}{m-1}\frac{z^n}{\rho^n}.
\end{align*}
Hence
\begin{equation}\label{eq:general-pole-coeff}
 \coeff{z}{n}A(z-\rho)^{-m}
 =A(-1)^m\binom{n+m-1}{m-1}\rho^{-n-m}.
\end{equation}
For a simple pole, \(m=1\), this reduces to
\(
 -A\rho^{-n-1}
\).  Since the residue of \((2-\e^z)^{-1}\) is \(A=-1/2\), its contribution is
\(
 (2\rho^{n+1})^{-1}
\).  Formula \eqref{eq:general-pole-coeff} also proves the polynomial factor in
\cref{thm:higher-poles}.

\section{First values and consistency checks}

\begin{center}
\begin{tabular}{@{}r r@{\qquad}r r@{}}
\toprule
\(n\)&\(\Fub_n\)&\(n\)&\(\Fub_n\)\\
\midrule
0&1&8&545835\\
1&1&9&7087261\\
2&3&10&102247563\\
3&13&11&1622632573\\
4&75&12&28091567595\\
5&541&13&526858348381\\
6&4683&14&10641342970443\\
7&47293&15&230283190977853\\
\bottomrule
\end{tabular}
\end{center}

All displayed values satisfy the following independent checks:
\begin{enumerate}
\item the finite Stirling sum \eqref{eq:def-fubini};
\item the initial condition and recurrence \eqref{eq:basic-rec};
\item the Taylor coefficients of \((2-\e^z)^{-1}\);
\item the geometric moment sum \eqref{eq:dobinski-fubini}.
\end{enumerate}
For every displayed index \(n\ge1\), they also satisfy the pole sum
\eqref{eq:exact-poles}, with the numerical truncation error bounded by
\eqref{eq:tail-hurwitz}.
The redundancy is useful in symbolic and numerical implementations because it
separates combinatorial indexing errors from analytic truncation errors.

\section*{Acknowledgment of source context}
\addcontentsline{toc}{section}{Acknowledgment of source context}
The combinatorial starting point and notation were aligned with the linked
\emph{Combinatorial Coefficient Calculus} in the ProveIt repository, whose
ordered-Bell section proves \eqref{eq:def-fubini}, \eqref{eq:basic-egf}, and
\eqref{eq:basic-rec}.  The organization of algebraic versus exponential scales
was informed by the repository's series-and-transseries drafts.  The pole
transseries, coefficient theorems, tail bounds, nonlinear multi-index calculus,
Borel dictionary, and higher-order-pole extension are developed here in a
self-contained form.

\begin{thebibliography}{99}

\bibitem{Bailey}
R.~W. Bailey,
\newblock The number of weak orderings of a finite set,
\newblock \emph{Social Choice and Welfare} \textbf{15} (1998), 559--562.
\newblock \url{https://doi.org/10.1007/s003550050123}.

\bibitem{Barthelemy}
J.-P. Barth\'elemy,
\newblock An asymptotic equivalent for the number of total preorders on a finite set,
\newblock \emph{Discrete Mathematics} \textbf{29} (1980), 311--313.
\newblock \url{https://doi.org/10.1016/0012-365X(80)90159-4}.

\bibitem{Comtet}
L. Comtet,
\newblock \emph{Advanced Combinatorics: The Art of Finite and Infinite Expansions},
\newblock D. Reidel, Dordrecht, 1974.

\bibitem{Costin}
O. Costin,
\newblock \emph{Asymptotics and Borel Summability},
\newblock Chapman \& Hall/CRC, 2009.

\bibitem{DLMF}
NIST Digital Library of Mathematical Functions,
\newblock Chapter 5, especially \S5.11, Asymptotic Expansions of the Gamma Function.
\newblock \url{https://dlmf.nist.gov/5.11}.

\bibitem{FlajoletSedgewick}
P. Flajolet and R. Sedgewick,
\newblock \emph{Analytic Combinatorics},
\newblock Cambridge University Press, 2009.

\bibitem{OEIS}
OEIS Foundation,
\newblock Entry A000670, Fubini numbers or ordered Bell numbers.
\newblock \url{https://oeis.org/A000670}.

\bibitem{Pippenger}
N. Pippenger,
\newblock The hypercube of resistors, asymptotic expansions, and preferential arrangements,
\newblock \emph{Mathematics Magazine} \textbf{83} (2010), 331--346.
\newblock Preprint: \url{https://arxiv.org/abs/0904.1757}.

\bibitem{ReshetnikovCCC}
V. Reshetnikov,
\newblock \emph{Combinatorial Coefficient Calculus: A Unified Formulae Monograph},
\newblock ProveIt repository, semi-formalized research-frontier draft, 2026.
\newblock \url{https://github.com/VladimirReshetnikov/ProveIt/tree/main/Analysis/FabiusFunction/docs/semi-formalized-research-frontiers/drafts/combinatorial-coefficient-calculus/Combinatorial_Coefficient_Calculus}.

\bibitem{ReshetnikovTransseries}
V. Reshetnikov,
\newblock \emph{Series and Transseries},
\newblock ProveIt repository, collection of semi-formalized research-frontier drafts, 2026.
\newblock \url{https://github.com/VladimirReshetnikov/ProveIt/tree/main/Analysis/FabiusFunction/docs/semi-formalized-research-frontiers/drafts/series-and-transseries}.

\bibitem{Sauzin}
D. Sauzin,
\newblock Introduction to 1-summability and resurgence,
\newblock arXiv:1405.0356, 2014.
\newblock \url{https://arxiv.org/abs/1405.0356}.

\bibitem{WhittakerWatson}
E.~T. Whittaker and G.~N. Watson,
\newblock \emph{A Course of Modern Analysis}, 4th ed.,
\newblock Cambridge University Press, 1927.

\end{thebibliography}

\end{document}
